%% ============================================================
%%  conclusion.tex — الخاتمة وملاحق ولائحة المراجع
%% ============================================================

%% ════════════════════════════════════════════════════════════
%%  الخاتمة
%% ════════════════════════════════════════════════════════════
\chapter*{الخاتمة}
\addcontentsline{toc}{chapter}{الخاتمة}
\markboth{الخاتمة}{الخاتمة}

خلص هذا التقرير، بعد استعراض شامل لمختلف جوانب الشغل المؤقت والمناولة في ضوء
مدونة الشغل المغربية، إلى جملة من النتائج الجوهرية:

\begin{enumerate}[leftmargin=*, itemsep=10pt]
    \item يُتيح النظام القانوني المغربي اللجوءَ إلى الشغل المؤقت في حالات محددة
          وحصرية، غير أن التطبيق الفعلي يكشف عن فجوات بين النص والممارسة.
    \item تُكرّس مدونة الشغل مبدأ التضامن بين شركة التشغيل المؤقت والمقاولة
          المستعملة، بوصفه ضمانة أساسية لحماية الأجراء.
    \item يظل دور مفتش الشغل محورياً، غير أنه مقيَّد بمحدودية الموارد البشرية
          والمادية.
    \item تُسهم التجارب المقارنة في إثراء النموذج المغربي وتجويده.
    \item تبقى التوصيات الموجّهة رهينةَ الإرادة السياسية والوعي الاجتماعي
          لدى مختلف الأطراف المعنية.
\end{enumerate}

\vspace{1cm}

إن تحقيق التوازن بين متطلبات المرونة الاقتصادية التي تنشدها المقاولات وضرورة
الحماية الاجتماعية التي تستحقها الطبقة العاملة، لَيُعدّ تحدياً حقيقياً يستوجب
حواراً اجتماعياً جاداً وإصلاحات تشريعية مستمرة.

%% ════════════════════════════════════════════════════════════
%%  الملاحق
%% ════════════════════════════════════════════════════════════
\appendix

\chapter{الملحق الأول: نصوص قانونية مختارة}
\addcontentsline{toc}{chapter}{الملحق الأول: نصوص قانونية مختارة}

%% أدرج هنا نصوص المواد القانونية ذات الصلة

\section{المادة 495 من مدونة الشغل}
\textit{(أدرج نص المادة كاملاً)}

\section{المادة 496 من مدونة الشغل}
\textit{(أدرج نص المادة كاملاً)}

\section{المادة 86 من مدونة الشغل}
\textit{(أدرج نص المادة كاملاً)}

%% ─────────────────────────────────────────────────────────
\chapter{الملحق الثاني: نموذج عقد التشغيل المؤقت}
\addcontentsline{toc}{chapter}{الملحق الثاني: نموذج عقد التشغيل المؤقت}

\textit{(أدرج هنا نموذج العقد المرجعي المستخدم في الوضعية المدروسة)}

%% ════════════════════════════════════════════════════════════
%%  لائحة المراجع
%% ════════════════════════════════════════════════════════════
\chapter*{لائحة المراجع}
\addcontentsline{toc}{chapter}{لائحة المراجع}
\markboth{لائحة المراجع}{لائحة المراجع}

%% ─── أولاً: النصوص القانونية ──────────────────────────────
\section*{أولاً: النصوص القانونية}
\begin{itemize}[leftmargin=*, itemsep=6pt]
    \item القانون رقم \textenglish{65.99} المتعلق بمدونة الشغل، الجريدة الرسمية،
          \textenglish{2003}.
    \item الظهير الشريف رقم \textenglish{1.03.194} بتاريخ \textenglish{14} رجب
          \textenglish{1424} (\textenglish{11} غشت \textenglish{2003}).
    \item المراسيم التطبيقية ذات الصلة.
\end{itemize}

%% ─── ثانياً: الكتب والمؤلفات ─────────────────────────────
\section*{ثانياً: الكتب والمؤلفات}
\begin{itemize}[leftmargin=*, itemsep=6pt]
    \item المؤلف، \textit{عنوان الكتاب}، دار النشر، المدينة، السنة.
    \item المؤلف، \textit{عنوان الكتاب}، دار النشر، المدينة، السنة.
\end{itemize}

%% ─── ثالثاً: المقالات والدراسات ──────────────────────────
\section*{ثالثاً: المقالات والدراسات}
\begin{itemize}[leftmargin=*, itemsep=6pt]
    \item المؤلف، «عنوان المقال»، \textit{اسم المجلة}، العدد، السنة، الصفحات.
    \item المؤلف، «عنوان المقال»، \textit{اسم المجلة}، العدد، السنة، الصفحات.
\end{itemize}

%% ─── رابعاً: المراجع الإلكترونية ─────────────────────────
\section*{رابعاً: المراجع الإلكترونية}
\begin{itemize}[leftmargin=*, itemsep=6pt]
    \item الموقع الرسمي لوزارة الإدماج الاقتصادي والمقاولة الصغيرة والتشغيل
          والكفاءات: \url{www.emploi.gov.ma}
    \item موقع منظمة العمل الدولية: \url{www.ilo.org}
\end{itemize}

